\subsection*{从物质结构到宏观电荷}

电与磁最初作为不同现象被研究，电磁感应的发现揭示了二者之间的联系，Maxwell 方程组又把
电场、磁场和光统一在同一理论中。静电学只处理电荷分布不随时间改变的情形，却是理解完整电磁
理论的起点。宏观物体是否带电，最终由内部正、负电荷的数量差决定；物体的尺寸和质量远大于
电子、质子等微观粒子，宏观电量却仍表现为元电荷的整数倍。

物质具有层级结构。原子由原子核和核外电子组成，原子核又由质子、中子构成，核子内部还存在
夸克与胶子等更深层结构。不同层级的典型长度相差许多数量级。讨论厘米尺度带电体时，不必追踪
每一个微观粒子，只需用总电荷、线电荷密度、面电荷密度或体电荷密度描述；讨论原子内部时，
连续分布近似便可能失效。图29.1(a)给出的长度标度用来说明模型选择，而不是把微观对象按实际比例
塞在同一张图中。

电荷守恒比“电子不会消失”更一般。成对产生、湮灭、核反应或粒子转化可以改变带电粒子的种类
和数目，但孤立系统的总电荷不变。电荷可加性意味着复合系统电量是各部分代数和；量子化则说明
通常可写 $Q=Ne$。在夸克层次会出现分数电荷，但自由宏观物体的净电荷仍以元电荷为基本尺度。

\subsection*{Coulomb 定律的适用范围与场的引入}

Coulomb 定律最初由扭秤实验定量建立。横杆绕悬丝转动，外部带电小球对杆端小球施加近似切向
的静电力，力矩 $Fd$ 与悬丝回复力矩 $\kappa\theta$ 平衡。如果把外电荷画在杆端小球的径向方向，
并把力也画成沿杆方向，力矩便会变为零，实验原理就被破坏。因此装置图中的电荷位置、力方向和
力臂必须同时满足几何关系。

点电荷模型要求带电体尺寸远小于相互距离。对球对称电荷分布，球外场可以等效为球心点电荷；
对一般形状则只能在足够远处作近似。Coulomb 定律还是静电规律：当电荷高速运动或分布随时间
变化时，磁效应和传播延迟不能忽略，不能继续把相互作用理解成瞬时的平方反比力。

引入电场以后，相互作用被分成两步：场源电荷建立空间中的场，试探电荷在当地电场中受力。场是
空间位置和时间的函数，带电粒子受到的力只由它所在位置的场决定。静止点电荷的场与 Coulomb
定律完全一致；对于变化电荷分布，场的描述仍然成立，而直接的“超距作用”图像不再可靠。

电场线只编码场的方向和相对强弱。静电场线不能相交、不能在无电荷处突然中断，也不能形成闭合
回路。振荡电荷产生的是随时间传播的电磁场，某一时刻的线形会随时间改变；因此原讲义把振荡
电荷与静态同号、异号电荷并列，是为了提醒读者电场概念并不限于静电场，而不是把三幅图都当作
时间不变的静电场线。

\subsection*{Gauss 定律的积分与微分图像}

闭合面的电通量是所有面积元上 $\vect E\cdot\dd\vect S$ 的代数和。面积元法向规定向外，场线
穿出给正贡献，穿入给负贡献。面外电荷的场线若穿过闭合面，必然一处进入、另一处离开，净贡献
为零；面内正电荷对应净流出，面内负电荷对应净流入。因此
\[
  \oint_S\vect E\cdot\dd\vect S=\frac{Q_{\rm in}}{\varepsilon_0}
\]
只取决于面内净电荷，与闭合面的形状以及面外电荷的位置无关。

高斯定律永远成立，但不总能方便地求场。只有球对称、无限柱对称或无限平面对称等情形，才能
在所选高斯面上把 $E$ 从积分中提出，或确认某些面上的通量为零。把“高斯定律成立”与“可用
高斯定律直接求 $E$”混为一谈，是静电题最常见的逻辑错误之一。

微分形式可以由一个边长分别为 $\Delta x,\Delta y,\Delta z$ 的小长方体理解。垂直于 $x$ 轴的
两个面，净通量在一阶近似下为
\[
  \frac{\partial E_x}{\partial x}\Delta x\Delta y\Delta z.
\]
另外两对表面分别给出 $\partial E_y/\partial y$ 和 $\partial E_z/\partial z$ 的贡献。六面相加，
再除以体积并令尺寸趋于零，便得到
\[
  \nabla\cdot\vect E=\frac{\rho}{\varepsilon_0}.
\]
图29.5(a)的立方体不仅是装饰，它明确了相对表面、外法向和差分位置；缺少这些位置关系，读者很难
从积分定律真正过渡到局部微分定律。

微分方程给出的是某一点附近电场的“源强”，并不能单独决定整个空间的场。相同的
$\nabla\cdot\vect E=\rho/\varepsilon_0$，配上不同边界条件会得到不同解；静电场还同时满足
$\nabla\times\vect E=0$。因此求实际电场时，需要把局部方程、电荷分布、导体表面条件和无穷远
条件放在一起。积分形式适合利用整体对称性，微分形式适合研究局部变化，二者是同一定律的不同
表达，而不是两条互不相干的公式。

跨过面电荷时，电场法向分量发生跃变
$E_{2n}-E_{1n}=\sigma/\varepsilon_0$，切向分量在静电条件下连续。这个边界关系可以由薄柱
高斯面直接得到，也是检查导体、带电薄膜和分区电场图是否画对的快速方法。

\subsection*{叠加原理与模型边界}

电场方程对场强是线性的，因此多个电荷分布产生的总场等于各部分单独产生的场之和。利用这一点，
复杂分布常被拆成“完整对称分布加上负的被挖去部分”，球形空腔、非均匀带电体和有限厚板都可
这样处理。叠加的是电场矢量；场强大小一般不能直接相加，必须先统一坐标分量。

“点电荷”“无限长直线”“无限大平面”都是带适用范围的模型。离带电体足够远时可用点电荷近似；
到直线中段的距离远小于线长时可忽略端点；观察尺度远小于平面尺寸时才可使用无限平面结果。若
题目进入端部、边缘或近场区域，对称性被破坏，高斯定律仍然成立，但不能再把场强当作高斯面上的
常量。先判断尺度关系，再选择模型，是使用高斯定律之前不可省略的一步。
